# Star catalog position and proper motion corrections in asteroid astrometry II: The Gaia era

# Abstract

Astrometric positions of moving objects in the Solar System have been measured using a variety of star catalogs in the past. Previous work has shown that systematic errors in star catalogs can affect the accuracy of astrometric observations. That, in turn, can influence the resulting orbit fits for minor planets. 

过去，天文学家曾使用多种星表来测量太阳系内移动物体的天体位置 。先前的研究表明，星表中的系统性误差会影响天体测量观测的准确性 ，这进而会影响到小行星的轨道拟合结果 。

In order to quantify these systematic errors, we compare the positions and proper motion of stellar sources in the most utilized star catalogs to the second release of the Gaia star catalog. The accuracy of Gaia astrometry allows us to unambiguously identify local biases and derive a scheme that can be used to correct past astrometric observations of solar system objects. 

为了量化这些系统性误差，我们将最常用的星表中恒星源的位置和自行（proper motion）与盖亚（Gaia）第二版星表进行比较 。盖亚天体测量学的高精度使我们能够明确地识别出局部偏差，并推导出一套方案，用以修正过去对太阳系天体的天体测量观测 。

Here we provide a substantially improved debiasing scheme for 26 astrometric catalogs that were extensively used in minor planet astrometry. Revised corrections near the galactic center eliminate artifacts that could be traced back to reference catalogs used in previous debiasing schemes. Median differences in stellar positions between catalogs now tend to be on the order of several tens of milliarcseconds (mas) but can be as large as 175 mas. Median stellar proper motion corrections scatter around 0.3 mas/yr and range from 1–4 mas/yr for star catalogs with and without proper motion, respectively. 

本文为26个在小行星天体测量中被广泛使用的星表提供了一套经过大幅改进的去偏差方案 。针对银心附近区域的修正消除了以往去偏差方案中因参考星表而引入的人为效应（artifacts） 。修正后，不同星表间恒星位置差异的中位数通常在几十毫角秒（mas）的量级，但最大可达175毫角秒 。恒星自行的修正中位数在 0.3 mas/yr 附近浮动；对于分别包含和不包含自行数据的星表，该修正范围在 1–4 mas/yr 之间 。

The tables in this work are meant to be applied to existing optical observations. They are not intended to correct new astrometric measurments as those should make use of the Gaia astrometric catalog. Since previous debiasing schemes already reduced systematics in past observations to a large extent, corrections beyond the current work may not be needed in the foreseeable future.

本研究提供的数据表旨在应用于已有的光学观测数据 。这些数据表不应用于修正新的天体测量，因为新的测量应当直接使用盖亚星表进行归算 。鉴于先前的去偏差方案已在很大程度上消除了历史观测数据中的系统性误差，因此在可预见的未来可能不再需要进行超出本工作范围的进一步修正 。

Keywords: Asteroids — Asteroids, dynamics — Orbit determination — Astrometry — Near- Earth objects

# 1. Introduction

Over more than two centuries astronomers have collected approximately 200 million astrometric observations of minor planets in the Solar System<sup>2</sup>. This dataset of astrometric measurements is essential to estimate the trajectories of Solar System Objects (SSOs) such as near- Earth asteroids (Milani and Grouchi, 2010). The vast majority of astrometry used for orbit determination of SSOs is optical, i.e. pairs of Right Ascension (RA) and Declination (DEC) measurements at given epochs. In contrast, only a tiny fraction of the SSO population has been observed by radar<sup>3</sup>. Improving the quality of optical astrometry is a worthwhile endeavor, especially if that means extending the arc of observations for objects of interest. 

在过去两个多世纪里，天文学家们已经收集了太阳系中约2亿次小行星的天体测量观测数据 。这个天体测量数据集对于估算近地小行星等太阳系天体（SSOs）的运行轨道至关重要 。用于SSOs轨道确定的绝大多数天体测量数据都是光学的，即在特定历元下的赤经（RA）和赤纬（DEC）成对测量值 。相比之下，只有极小一部分的SSO曾被雷达观测过 。因此，提高光学天体测量的质量是一项非常有价值的工作，特别是当这意味着可以延长对感兴趣天体的观测弧长时 

Over the years a variety of star catalogs has been used to measure astrometric positions of asteroids to a common frame of reference such as the ICRF (Ma et al., 2009) (Table 1). With improving quality of astrometric measurements it became evident that some observations of asteroids showed systematic errors that deviated from the expected scatter around the nominal trajectory (Carpino et al., 2003). Some of those systematic errors could be traced back to the use of specific astrometric catalogs in the observation reduction process. To assess these systematic errors one can compare stellar positions and proper motion between various star catalogs. Chesley et al. (2010) used the 2MASS (Skrutskie et al., 2006) catalog as a reference to identify local systematics in positions of stars in the USNO catalogs (Monet, 1998; Lasker et al., 1996a; Monet et al., 2003). Correcting for regional biases in star positions led to better ephemeris predictions and improved the error statistics of astrometric observations of asteroids.

多年来，天文学家们使用多种星表，将小行星的天体位置测量归算到一个共同的参考框架下，例如国际天球参考框架（ICRF） 。随着测量质量的提升，一些小行星的观测数据明显表现出系统性误差，其分布偏离了理论轨道周围的预期散布范围 。其中一些系统性误差可以追溯到在观测归算过程中使用了特定的星表 。评估这些系统性误差的方法之一是比较不同星表之间的恒星位置和自行 。Chesley等人（2010）使用2MASS星表作为参考，识别出了USNO系列星表中恒星位置的局部系统性偏差 。修正恒星位置的区域偏差可以得到更好的星历预测，并改善小行星天文观测的误差统计。

Selecting the most accurate stars in the PPM XL catalog based on 2MASS astrometry, Farnocchia et al. (2015) improved the scheme by Chesley et al. (2010) and debiased practically all of the prevalent star catalogs used for asteroid astrometry. Farnocchia et al. (2015) included proper motion in the debiasing process which led to widespread improvements in the post- fit residuals of asteroids, most notably (99942) Apophis, (101955) Bennu and (6489) Golevka. The present work follows a similar approach with the advantage of having Gaia data, described in section 2, as a reference. Following a brief outline of the debiasing process (section 3) we calculate position and proper motion corrections for stellar catalogs that have been used extensively in minor planet astrometry (sections 4 - 6). In section 7 we then assess the quality of ephemeris predictions based on the new debiasing scheme. A discussion of our results concludes this article.

Farnocchia等人（2015）通过在2MASS天体测量的基础上筛选PPM XL星表中精度最高的恒星，改进了Chesley等人（2010）的方案，并对几乎所有用于小行星天体测量的主流星表进行了去偏差处理 。Farnocchia等人（2015）将自行纳入去偏差过程，这使得小行星的拟合后残差得到了广泛改善，其中最为显著的是(99942)阿波菲斯、(101955)贝努和(6489)格列佛卡 。本研究遵循了类似的方法，但优势在于使用第2节中所述的盖亚（Gaia）数据作为参考 。在简要概述去偏差过程（第3节）之后，我们为那些在小行星天体测量中被广泛使用的星表计算了位置和自行修正（第4-6节） 。接着在第7节中，我们评估了基于新去偏差方案的星历预报质量 。最后，我们通过对结果的讨论来为本文作结 。

# 2. The Gaia catalogs

Since 2013 the Gaia mission has been collecting astrometric information on roughly 1.7 billion sources (Gaia Collaboration et al., 2018). A first data release based on a partially completed survey featured around a billion sources largely without proper motion (Lindegren et al., 2016). With the second Gaia data release (Gaia DR 2) published on April 25, 2018 this information has become publicly accessible through the Gaia Archive<sup>4</sup>. About 1.3 billion sources between 3<Gmag<21 come with a full five-parameter astrometric solution, i.e. right ascension and declination positions on the sky  $(\alpha , \delta)$ , parallaxes, and proper motion. Uncertainties in the parallaxes range between 0.04 milliarcseconds (mas) for sources with Gmag<15 and 0.7 mas at Gmag=20. The corresponding uncertainties in the proper motion components start at 0.06 mas/yr for Gmag<15 and reach 1.2 mas/yr for Gmag=20. The Gaia astrometric catalog is, thus, a reference of unprecedented accuracy well suited to identify systematic errors in other star catalogs.

自2013年以来，盖亚（Gaia）任务一直在收集约17亿个天体源的天体测量信息 。其初版数据发布基于一次未完成的巡天，包含了约10亿个天体源，但大部分没有自行（proper motion）数据 。随着2018年4月25日盖亚第二版数据（Gaia DR2）的发布，这些信息已通过盖亚档案馆（Gaia Archive）向公众开放 。其中，约13亿个G星等在3至21等之间的天体源具备完整的五参数天体测量解，即天球上的赤经（α）和赤纬（δ）位置、视差（parallax）以及自行 。对于G星等小于15的天体源，其视差不确定度在0.04毫角秒（mas）左右；对于G星等为20的天体源，不确定度则为0.7毫角秒 。相应的，自行分量的不确定度对于G星等小于15的天体源始于0.06 mas/yr，到G星等为20时则达到1.2 mas/yr 。因此，盖亚星表是一个精度空前的参考基准，非常适合用于识别其他星表中的系统性误差 

# 3. Star position and proper motion corrections

Similar to Chesley et al. (2010) and Farnocchia et al. (2015) we compare the various star catalogs that were used for minor planet astrometry in the past to a reference catalog, in this case Gaia DR 2. By dividing the celestial sphere into 49,152 equal- area tiles ( $\sim$ 0.8 square degrees each) using a HEALPix tesselation (Gorski et al., 2005) and comparing stellar positions and proper motion within a given tile we can quantify local systematic differences. 

类似于 Chesley 等（2010）和 Farnocchia 等（2015）的研究，我们将过去用于小行星天体测量的各类恒星目录与一个参考目录进行比较，本研究采用的参考目录为 Gaia DR2。具体方法是利用 HEALPix 镶嵌（Gorski et al., 2005）将全天球划分为 49,152 个等面积分区（每个分区约为 0.8 平方度），并在各个分区内比较恒星位置与自行，从而定量评估局部系统性差异。

To this end we extract stellar positions and proper motion at epoch J2000 from the reference and the test catalog. Given the variability in the number of stars per tile in each catalog, identification of common stars requires some attention. We first search for spatial correlation between astrometric sources within 1" of the reference position. If candidates are found we apply the identification and outlier rejection criterion described in Farnocchia et al. (2015) to ensure spurious matches are excluded from the sample. The median of the differences in the astrometric quantities of each identified star per tile yields the local corrections for position ( $\Delta \alpha_{J2000}, \Delta \delta_{J2000}$ ), and proper motion ( $\Delta \mu_{\alpha}, \Delta \mu_{\delta}$ ) at epoch I2000. 

为此，我们从参考目录与测试目录中提取 历元 J2000 的恒星位置和自行。鉴于各目录在每个分区中恒星数量的差异，公共恒星的识别需要特别注意。我们首先在参考位置 1″ 范围内搜索天体测量源的空间相关性；若发现候选者，则应用 Farnocchia 等（2015）所描述的识别与离群值剔除准则，以确保排除伪匹配。随后，在每个分区内取已识别恒星的天体测量量差值的中位数，即可得到该分区在历元 J2000 的位置修正量（$\Delta \alpha_{J2000}, \Delta \delta_{J2000}$）以及自行修正量（$\Delta \mu_{\alpha}, \Delta \mu_{\delta}$）。

Astrometric observations of minor planets can then be corrected for catalog systematics by subtracting

小行星的天体测量观测值即可通过减去下列量来修正因目录系统性误差所带来的偏差：
$$
\Delta \alpha &= \Delta \alpha_{J2000} + \Delta \mu_{\alpha}\cdot \Delta t \\
\Delta \delta &= \Delta \delta_{J2000} + \Delta \mu_{\delta}\cdot \Delta t
$$

with  $\Delta t = t - 2451545.0$ , the difference between the epoch of observation and J2000 in JD. All differential quantities  $\Delta \alpha$ ,  $\Delta \alpha_{J2000}$ , and  $\Delta \mu_{\alpha}$  include the spherical metric factor  $\cos \delta$ .

其中，$\Delta t = t - 2451545.0$，即观测历元与 J2000 历元之间的儒略日差值。所有差分量 $\Delta \alpha$、$\Delta \alpha_{J2000}$ 和 $\Delta \mu_{\alpha}$ 均包含球面度量因子 $\cos \delta$。

# 4. Added astrometric catalogs

In addition to the 20 catalogs discussed in Farnocchia et al. (2015) and Gaia DR 1 we added the following five catalogs to be compared to Gaia DR 2:

除 Farnocchia 等（2015）讨论的 20 个星表及 Gaia DR1 外，我们在本研究中还增加了以下五个需要与 Gaia DR2 进行比较的星表：

- CMC 15 is the  $15^{\text{th}}$  installment of the Carlsberg Meridian Catalogue generated through the Carlsberg Meridian Telescope, formerly the Carlsberg Automatic Meridian Circle (Copenhagen University et al., 2011). CMC 15 has been compiled between March 1999 and March 2011, encompasses around 134,000,000 entries, and claims a positional accuracy between 35 to 100 mas for a declination range from -40 and  $+50$  degrees.
- CMC 15：即第 15 版 Carlsberg 子午环目录，由 Carlsberg 子午望远镜（前称 Carlsberg 自动子午环）生成（Copenhagen University et al., 2011）。CMC 15 编纂于 1999 年 3 月至 2011 年 3 月之间，包含约 1.34 亿条目，在赤纬范围 $-40^\circ$ 至 $+50^\circ$ 内，其位置精度声称在 35 至 100 mas 之间。
- The  $8^{\text{th}}$  data release of the Sloan Digital Sky Survey (SDSS) added roughly 4 times the number of sources published in data release 7, namely about 325 million (Adelman-McCarthy and et al., 2011). The fact that astrometric data together with proper motion is available has made SDSS DR 7 (Abazajian et al., 2009) attractive for minor planet astrometry. Almost 1 million observations of minor planets were reduced with SDSS DR 7 (Table 1). SDSS DR 8, although offering a significant increase in source density, has not been as popular, which is partly due to known issues with the astrometry (Aihara et al., 2011).
- Sloan 数字巡天（SDSS）第 8 次数据发布（DR8）：相比于第 7 次数据发布，SDSS DR8 增加了约 4 倍的源，总数达到约 3.25 亿（Adelman-McCarthy et al., 2011）。由于同时提供了天体测量数据和自行信息，SDSS DR7（Abazajian et al., 2009）在小行星天体测量中曾非常有吸引力，约有近 100 万条小行星观测数据是基于 DR7 处理的（见表 1）。尽管 DR8 提供了更高的源密度，但其受欢迎程度却不及 DR7，部分原因在于其存在已知的天体测量问题（Aihara et al., 2011）。
- The Space Surveillance Telescope (SST) has been relying on a proprietary catalog based on selected sources from a variety of other astrometric and photometric catalogs for the reduction of its 14+ million observations of asteroids in the Solar System. Although not in the public domain, the  $5^{\text{th}}$  version of the SST catalog (SST-RC5) has been kindly made available for this study (Monet, 2018). Most of the astrometric observations considered in this work have been submitted under SST-RC4. Astrometric standard stars in SST- RC5 are largely the same as in the  $4^{\text{th}}$  installment of the SST catalog, however (Monet, 2018). Hence, we use the same correction to debias astrometric observations submitted under SST- RC4 and SST- RC5.
- 空间监视望远镜（SST）星表：SST 一直依赖于一个专有星表，该星表基于从多个天体测量及光度星表中精选的源构建，用于处理其超过 1,400 万条小行星观测数据。尽管该星表并未公开，但第 5 版 SST 星表（SST-RC5）已慷慨地提供给本研究使用（Monet, 2018）。本研究中考虑的大多数观测数据是以 SST-RC4 提交的。然而，SST-RC5 的天体测量标准星与第 4 版星表中的基本一致（Monet, 2018）。因此，我们对 SST-RC4 与 SST-RC5 提交的观测数据采用相同的改正方案进行去偏处理。
- UCAC 5, is the  $5^{\text{th}}$  installment of the US Naval Observatory CCD Astrograph Catalog (UCAC) (Zacharias et al., 2017). By combining UCAC 5 with Gaia DR1 data new proper motions on the Gaia coordinate system for over 107 million stars were obtained with typical accuracies of 1-2 mas/yr for R=11-15 mag, and about 5 mas/yr through R=16 mag.
- UCAC 5：即美国海军天文台 CCD 天体照相机星表（UCAC）的第 5 版（Zacharias et al., 2017）。通过将 UCAC 5 与 Gaia DR1 数据结合，在 Gaia 坐标系上获得了超过 1.07 亿颗恒星的新自行，其典型精度在 R=11–15 星等范围内为 1–2 mas/yr，而在 R=16 星等时约为 5 mas/yr。
- URAT, a follow-up project to the UCAC series featured and increased limiting magnitude that allowed for a roughly 4-fold increase in the average number of stars per square degree as compared to UCAC. Additionally, stars as bright as about  $3^{\text{rd}}$  magnitude were added to the catalog. The URAT 1 catalog (Zacharias et al., 2015) has its mean epoch between 2012.3 and 2014.6 supporting a magnitude range from 3 to 18.5 in R-band with a positional precision of 5 to 40 mas. It covers most of the northern hemisphere and some areas down to -24.8 deg in declination.
- URAT 1：作为 UCAC 系列的后续项目，URAT 扩展了星表的极限星等，使得单位平方度的平均恒星数相较于 UCAC 增加了约 4 倍。此外，其还补充了亮至约 3 等的恒星。URAT 1 星表（Zacharias et al., 2015）的平均历元介于 2012.3 与 2014.6 之间，覆盖 R 波段 3 至 18.5 等的恒星，位置精度在 5 至 40 mas 之间。其覆盖范围包括北半球大部分区域以及赤纬 $-24.8^\circ$ 以内的部分南天。

Table 2 shows that roughly 1% of observations do not have a catalog escriptor and can, therefore, not be corrected. About 300,000 observations have been submitted without specifying which version of the GSC catalog was used. Since the various installments of the GSCs have different correction tables, those observations are difficult to debias as well. Observations reduced by catalogs not explicitly mentioned in Table 1 account for less than 0.05% of all database entries.Their weight on the bulk of orbit solutions is deemed small enough not to merit explicit debiasing.

表 2 显示，约 1% 的观测数据缺乏星表说明，因此无法进行改正。此外，约有 30 万条观测数据未明确指出所使用的 GSC 星表版本。由于不同版本的 GSC 具有不同的改正表，这部分观测数据也难以进行去偏处理。至于那些由表 1 中未明确列出的星表所归算的观测数据，其比例不足数据库总条目的 0.05%。这部分观测对整体轨道解的影响被认为微不足道，因此不值得进行明确的去偏处理。

# 5. Catalog Comparison Results

With more than 40 million entries in the Minor Planet Center database the USNO A2.0 catalog has been one of the most popular choices for reducing astrometric observations of minor planets. Figures 1 and 2 show the proposed corrections derived from a comparison with Gaia DR 2. The regional biases in stellar positions first found by Chesley et al. (2010) and later confirmed by Farmocchia et al. (2015) are clearly visible in the top panels of Figure 1. Figure 2 shows the distribution of the corrections collected over the entire sky. The visual impression from Figure 1, namely that there is an overall bias towards higher declination values in stellar positions, is confirmed in the distribution of  $\Delta \delta_{J2000}$ . The lack of proper motion as well as a median bias of 162 mas in declination with a mode closer to a quarter arcsecond all suggest that astrometric observations reduced with USNO A2.0 could benefit from a correction. Table 2 reports the size of the corrections in terms of the all sky median (MED) and median absolute deviation (MAD) for each studied catalog. 

在小行星中心（MPC）数据库中，USNO A2.0 星表包含超过 4000 万条记录，是用于小行星天体测量归算最常用的星表之一。图 1 与图 2 展示了基于 Gaia DR2 对比所得的改正量。Chesley 等（2010）首次发现、并由 Farnocchia 等（2015）进一步确认的恒星位置区域性系统偏差，在图 1 的上半部分中清晰可见。图 2 显示了在全天范围内收集的改正分布。由图 1 所呈现的总体印象——即恒星位置在赤纬方向存在系统性正偏差——在 $\Delta \delta_{J2000}$ 的分布中得到了验证。缺乏自行信息、以及赤纬方向中位偏差达到 162 mas（众数接近 0.25″）的事实，都表明利用 USNO A2.0 归算的天体测量观测可通过改正获得显著改进。表 2 给出了每个研究星表在全天范围内的中位数（MED）及中位绝对偏差（MAD）的修正量。

Of particular interest is the comparison between PPM XL and Gaia DR 2, since a subset of the former was used as a reference for stellar positions and proper motion in Farnocchia et al. (2015). The corresponding graphs are presented in Figure 3. The range of position corrections is only about one sixth of that in USNO A2.0 and the subset of 2MASS stars used for astrometry also has a relatively small positional bias as can be gathered from Table 2. A comparison between Gaia DR 2 and UCAC 4 (Figure 5) reveals that the relatively large corrections seen by Farnocchia et al. (2015, Figures 9 to 12) near the galactic center ( $\alpha \approx 270$  deg,  $\delta \approx - 30$  deg) appear to actually be associated to the PPM XL reference, rather than the UCAC catalogs. This is indicated by the signature seen near the galactic center in all four panels of Figure 3 that is now substantially absent from all four panels of Figure 5.

特别值得注意的是 PPM XL 与 Gaia DR2 的比较，因为在 Farnocchia 等（2015）的研究中，PPM XL 的一个子集曾被用作恒星位置和自行的参考。相关图像见图 3。其位置改正的范围仅为 USNO A2.0 的约六分之一，而 2MASS 子集中用于天体测量的恒星在位置上的系统偏差也相对较小，这可由表 2 得出。Gaia DR2 与 UCAC 4 的对比（见图 5）显示，Farnocchia 等（2015，图 9–12）报告的在银河中心附近（$\alpha \approx 270^\circ, \delta \approx -30^\circ$）所见的较大改正，实际上应归因于 PPM XL 参考，而非 UCAC 星表。这一点由图 3 四幅面板在银河中心区域呈现的特征所表明，而在图 5 的四幅面板中，该特征几乎消失。

With the exception of the artificial structures close to the Galactic center introduced through PPM XL, our results for UCAC, USNO and CMC catalogs are very similar to those found in Farnocchia et al. (2015). The corresponding Figures are provided in the Appendix. 

除去由 PPM XL 引入的银河中心附近的人工结构之外，我们对 UCAC、USNO 及 CMC 星表的结果与 Farnocchia 等（2015）的结论非常一致。相关图像在附录中给出。

Among the newly added catalogs UCAC 5 continues the tradition of the UCAC family with an increase in the overall quality of the catalog. Local corrections are mostly due to differences between stellar positions in Gaia DR 2 and Gaia DR 1. The latter were ingested into UCAC 5. URAT 1 has a higher source density compared to UCAC 4, but it also exhibits larger deviations from Gaia DR 2 in both stellar positions and proper motion. We confirm the known issues with astrometric accuracy around the north celestial pole in SDSS DR 8 (Aihara et al., 2011), which were addressed in the following data releases. 

在新纳入的星表中，UCAC 5 延续了 UCAC 系列的传统，并提升了整体质量。其局部改正主要源自 Gaia DR2 与 Gaia DR1 恒星位置的差异，因为后者被纳入了 UCAC 5。URAT 1 相较于 UCAC 4 拥有更高的源密度，但其在恒星位置与自行上均表现出与 Gaia DR2 更大的偏差。我们也确认了 SDSS DR8 在北天极区域存在已知的天体测量精度问题（Aihara et al., 2011），这些问题已在后续数据发布中得到修正。

Apart from minor systematics stemming from survey incompleteness, Gaia DR 1 had excellent stellar positions at the epoch of its publication (Lindegren et al., 2016). Since Gaia DR 1 largely lacks proper motion, the differences to Gaia DR 2 become significant at the epoch of comparison (J2000), however. This is reflected in Table 2. The quality of stellar positions, even if measured through the Gaia satellite, degrades relatively quickly with time if no proper motion corrections are applied. Minor planet astrometrists are, thus, encouraged to switch to Gaia DR 2 as soon as possible. We debias all observations where corrections are available with the exception of those reduced with ACT, Gaia DR 1, Tycho- 2 and UCAC- 5, since the latter show no large scale structures in local position and proper motion systematics.

除去因巡天不完整性引起的轻微系统效应外，Gaia DR1 在其发布历元的恒星位置表现优异（Lindegren et al., 2016）。然而，由于 Gaia DR1 在很大程度上缺乏自行信息，其与 Gaia DR2 在历元 J2000 的对比差异显著，这一点反映在表 2 中。即使是由 Gaia 卫星测得的恒星位置，若未进行自行修正，其精度也会随时间迅速劣化。因此，建议小行星天体测量学者尽快切换至 Gaia DR2。我们对所有存在改正的观测数据进行了去偏处理，唯独 ACT、Gaia DR1、Tycho-2 及 UCAC 5 例外，因为这些星表在局部位置和自行系统误差中未表现出大尺度结构。

# 6. Smoothing of debiasing tables

Due to the discrete nature of calculating corrections to astrometric measurements on HEALPix tiles, transitions between neighboring corrections are not guaranteed to be smooth. As an example, Fig. 7 shows that adjacent USNO A2.0 tiles can have significantly different bias values. Corrections for minor planet observations that cross such a border would see an artificial jump in the corresponding coordinate. A finer HEALPix tessellation could be used to mitigate such a problem. However, a finer tessellation would mean that fewer stars per tile are available to estimate the magnitude of the correction. 

由于在 HEALPix 分区上计算天体测量改正具有离散性，相邻分区之间的改正量过渡未必平滑。举例而言，图 7 显示，相邻的 USNO A2.0 分区可能存在显著不同的偏差值。对于跨越这种边界的小行星观测，其对应坐标的改正可能会出现人为的跳变。采用更精细的 HEALPix 镶嵌可以缓解这一问题，但更高的分辨率也意味着每个分区可用于估计改正量的恒星数量减少。

In order to circumvent reducing the sample size of stars per tile, we use radial basis functions (RBFs, Schaback, 1995) to interpolate data from neighboring tiles onto a finer HEALPix tessellation. To this end, we increased the number of tiles from 49,152 to 786,432 and subsequently interpolated the catalog corrections from the coarser tessellation onto the new tile centers. This process reduces artificial discontinuities along tile borders while at the same time retaining a sufficient number of stars in each of the stencil tiles. 

为避免样本数量的减少，我们采用径向基函数（RBFs，Schaback, 1995）在更精细的 HEALPix 镶嵌上对相邻分区的数据进行插值。具体而言，我们将分区数从 49,152 增加至 786,432，并将粗分辨率镶嵌上的目录改正插值到新的分区中心。该过程在保证每个模板分区内恒星数量足够的同时，降低了分区边界处的人工不连续性。

In Fig. 8, the distribution of inter-tile corrections in RA is shown for USNO A2.0 before and after smoothing. The smoothing of catalog corrections leads to a substantial reduction of the variance (≈90%) and maximum magnitude (≈40%) in the distribution of discontinuous transitions between HEALPix tiles. The results for the DEC coordinate are essentially identical and therefore omitted. Reductions in the discontinuities in other catalogs due to smoothing are similar to those shown for USNO-A2.0.

图 8 展示了 USNO A2.0 在平滑前后赤经方向上跨分区改正的分布。对目录改正进行平滑处理后，不连续跃迁分布的方差减少了约 90%，最大幅度减少了约 40%。赤纬方向的结果基本相同，因此未单独列出。对于其他星表，通过平滑所减少的不连续性与 USNO A2.0 的结果相似。

# 7. Ephemeris Prediction Test

We follow the example of Chesley et al. (2010, sec. 6) and Farnocchia et al. (2015, sec. 4.3) and evaluate the impact of the new debiasing scheme on asteroid orbit determination and prediction. To achieve that we compare the prediction performance for an ensemble of 700 near- Earth objects (NEOs), each of which had at least five apparitions since 1998. Every one of those apparitions has at least ten observations over 15 days. Using the astrometric observations from either the first two apparitions or the penultimate apparition, respectively, we construct predictions close to the center of the third apparition, near the middle of the true observational arc. The predictions are then compared to the reference solution that uses the full arc. The analysis is performed in a consistent way such that a given debiasing scheme is used for both the short arcs derived from a subset of apparitions and the full arc. All astrometric data is weighted according to the scheme presented in Veres et al. (2017).

我们借鉴 Chesley 等（2010，第 6 节）与 Farnocchia 等（2015，第 4.3 节）的做法，评估新的去偏方案对小行星轨道确定与预测的影响。为此，我们对一个包含 700 颗近地天体（NEOs） 的样本进行预测性能比较。这些天体自 1998 年以来均至少出现过五次，每次出现均有不少于 15 天内的 10 次观测。我们分别利用其前两次出现或倒数第二次出现的天体测量观测，构建接近第三次出现中点（即实际观测弧段中点附近）的预测，并将预测结果与使用完整观测弧所得的参考解进行比较。该分析过程保持一致性，即对于由部分出现构建的短弧和完整弧均使用同一去偏方案。所有天体测量数据的加权采用 Vereš 等（2017）提出的方案。

Figure 8 shows the corresponding results. The prediction errors in the orbital semimajor axes of 700 NEOs derived from astrometry debiased with respect to Gaia DR 2 are only slightly smaller that those using the debiasing scheme presented in Farnocchia et al. (2015). These results suggest that most of the statistical bias was de facto eliminated with the previous debiasing approach. Using interpolated debiasing tables as discussed in the previous section did not significantly affect the above prediction statistics, either. Given the larger size of the interpolated table and the potential loss in computational efficiency we recommend using interpolated values only when quasi- discontinuities in the former debiasing table are an obstacle to high- fidelity orbit determination.

图 8 展示了相应结果。基于 Gaia DR2 去偏后的天体测量所得的 700 颗 NEO 轨道半长轴预测误差，仅比 Farnocchia 等（2015）所提出的去偏方案下的结果略小。这表明，大多数统计偏差事实上已经通过之前的去偏方法消除。使用上一节所述的插值去偏表，对上述预测统计量的影响也并不显著。考虑到插值表体量较大且可能导致计算效率降低，我们建议仅在原去偏表的准不连续性影响高精度轨道确定时，才使用插值值。

# 8. Conclusions

The Gaia astrometric catalog published with the second Gaia data release has become the new standard for minor planet astrometry. Its uniformity and high quality in both stellar positions and proper motion allow us to improve prior astrometric observations of minor planets by correcting for potential systematics in other astrometric catalogs. In this work we have used Gaia DR 2 as a reference to calculate debiasing tables for 26 astrometric catalogs. New tables for six catalogs (PPM XL, SST RC5, Gaia DR 1, CMC 15, URAT, UCAC 5 and SDSS 8) that had not been studied in Farnocchia et al. (2015) are included in this work. The quality of catalog corrections presented here is significantly improved compared to previous work, in particular for the UCAC series. 

伴随 Gaia 第二次数据发布（Gaia DR2） 所公布的天体测量星表，已成为小行星天体测量的新标准。其在恒星位置与自行上的均一性与高质量，使我们能够通过修正其他天体测量星表中潜在的系统性误差，改进既有的小行星天体测量观测。在本研究中，我们使用 Gaia DR2 作为参考，为 26 个天体测量星表计算了去偏表。其中包括 6 个在 Farnocchia 等（2015）中未被研究的星表（PPM XL、SST RC5、Gaia DR1、CMC 15、URAT、UCAC 5 以及 SDSS 8）。与此前工作相比，本研究所提供的目录改正质量显著提高，尤其是 UCAC 系列。

The switch to the new debiasing scheme based on Gaia DR 2 delivers only incremental improvement of the quality of near- Earth asteroid orbits over Farnocchia et al. (2015), however. This suggests that further debiasing attempts may not be necessary. The procedure suggested in this work can help mitigate some of the systematics introduced by the use of star catalogs other than Gaia. We would like to emphasize that the tables are meant as an a posteriori correction to existing astrometric data. They should not be used to debias new astrometric measurements. We strongly recommend contemporary Solar System astrometry be conducted with Gaia DR 2 and its successors.

然而，基于 Gaia DR2 的新去偏方案，在近地小行星轨道质量上的提升相较 Farnocchia 等（2015）仅为增量性。这表明进一步的去偏尝试可能已无必要。本研究所提出的程序有助于缓解因使用非 Gaia 星表而引入的部分系统误差。我们强调，这些去偏表旨在对既有的天体测量数据进行事后修正，而不应用于对新的天体测量进行去偏。我们强烈建议当代太阳系天体测量应直接采用 Gaia DR2 及其后续版本。

# 9.Data Availability

Tables containing standard resolution and interpolated catalog corrections derived in this work are publicly available at ftp://ssd.jpl.nasa.gov/pub/ssd/debias/debias_2018.tgz and ftp://ssd.jpl.nasa.gov/pub/ssd/debias/debias_hires2018.tgz.

本研究中推导的标准分辨率与插值目录改正表可通过以下公开地址获取：

- ftp://ssd.jpl.nasa.gov/pub/ssd/debias/debias_2018.tgz
- ftp://ssd.jpl.nasa.gov/pub/ssd/debias/debias_hires2018.tgz

# Acknowledgments

This research has received funding from the Jet Propulsion Laboratory through the California Institute of Technology postdoctoral fellowship program, under a contract with the National Aeronautics and Space Administration. VizieR catalogue access tools, CDS, Strasbourg, France, have been essential to this research project (Ochsenbein et al., 2000). This work made use of the IAU Minor Planet Center observations database as well as results from the European Space Agency (ESA) space mission Gaia. Gaia data are being processed by the Gaia Data Processing and Analysis Consortium (DPAC). Funding for the DPAC is provided by national institutions, in particular the institutions participating in the Gaia MultiLateral Agreement (MLA). The Gaia mission website is https://www.cosmos.esa.int/gaia. The Gaia archive website is https://archives.esac.esa.int/gaia.

本研究获得了喷气推进实验室（Jet Propulsion Laboratory）通过加州理工学院博士后奖学金项目的资助，并与美国国家航空航天局（NASA）签订合同。VizieR 星表访问工具（CDS, Strasbourg, France；Ochsenbein et al., 2000）在本研究中发挥了关键作用。本研究同时利用了国际天文学联合会（IAU）小行星中心的观测数据库，以及欧洲空间局（ESA）空间任务 Gaia 的成果。Gaia 数据由 Gaia 数据处理与分析联盟（DPAC） 处理，DPAC 的资金由各国机构提供，尤其是参与 Gaia 多边协议（MLA） 的机构。

Gaia 任务网站：https://www.cosmos.esa.int/gaia
 Gaia 数据档案网站：https://archives.esac.esa.int/gaia

# Appendix

The appendix shows local position and proper motion systematics for stellar catalogs not presented in the main body of the article (Figures 9- 24). For catalogs without proper motion only the bias values in right ascension and declination are provided. Proper motion corrections for those catalogs resemble those of USNO A2.0 displayed in Figure 1.

附录展示了未在正文中呈现的恒星星表的局部位置与自行系统误差（见图 9–24）。对于缺乏自行信息的星表，仅提供赤经与赤纬方向的偏差值。此类星表的自行改正与 USNO A2.0 在图 1 中所示的结果相似。







WE
**ast**

hipparcos

**actrc**

sao

agk3

gaia dr3

gaia edr3

**ac**

**yale**

**poss**

**acrs**

**agk1**

positions are corrected for full relativistic aberration but not for relativistic light deflection in the gravitational field of the Solar System



EGGL

tycho2

act

ucac4

ucac2

gaia dr1

usnoa2.0

gsc2.2

ppm